% ECONOMICS_PROBLEM: {"id": "O31-510", "title": "Entrant versus incumbent innovation", "jel_code": "O31", "source_id": "OP-0266"}
% !TeX program = xelatex
\documentclass[11pt,a4paper]{article}
\usepackage[margin=23mm]{geometry}
\usepackage{fontspec}
\setmainfont{Latin Modern Roman}
\usepackage{amsmath,amssymb,mathtools}
\usepackage{xcolor,enumitem,booktabs,longtable,array,xurl,hyperref}
\definecolor{muted}{HTML}{53636C}
\hypersetup{colorlinks=true,urlcolor=blue,linkcolor=blue}
\setlength{\parindent}{0pt}
\setlength{\parskip}{5pt}
\newcommand{\R}{\mathbb R}
\newcommand{\E}{\mathbb E}
\newcommand{\N}{\mathbb N}
\newcommand{\ind}{\mathbf 1}
\newcommand{\Var}{\operatorname{Var}}
\newcommand{\Cov}{\operatorname{Cov}}
\newcommand{\supp}{\operatorname{supp}}
\DeclareMathOperator*{\argmin}{arg\,min}
\DeclareMathOperator*{\argmax}{arg\,max}
\newcommand{\entrymeta}[3]{{\sffamily\small\color{muted}\textbf{\texttt{#1}}\quad|\quad #2\par #3\par}}
\newcommand{\Problemref}[1]{\texttt{#1}}
\newcommand{\JELref}[2]{\texttt{#2} (JEL \texttt{#1})}
\begin{document}
\section*{Entrant versus incumbent innovation}
\textbf{Problem ID:} \texttt{O31-510}\quad \textbf{JEL:} \texttt{O31}\par
% BEGIN ECONOMICS STATEMENT
\label{problem:OP-0266}
{\sffamily\small\color{muted}Original subject group: Growth and productivity\par}
\label{definitions:problem:30007547}
\textbf{Audit classification:} Explicit published open question. Complementarity of creative destruction and incumbent innovation is explicit; the derivative design is editorial. Verification concerns the cited version, not first priority or proof that this exact editorial specification remains unsolved on October 8, 2026.
\subsection*{Definitions and precise research target}
\paragraph{Editorial mathematical specification.} Fix a product-line innovation model with incumbents choosing research intensity \(r\ge0\) and potential entrants generating creative-destruction hazard \(\lambda\ge0\). Let \(V(r,\lambda;w,p)\) be incumbent expected discounted profit net of research cost, with wage \(w\) and demand parameters \(p\) initially fixed. Under a unique differentiable interior best response \(r^*(\lambda;w,p)\), define strategic complementarity or substitutability by the sign of
\[
\frac{\partial r^*(\lambda;w,p)}{\partial\lambda}.
\]
Identify this derivative over a declared range of entrant pressure and distinguish it from the general-equilibrium derivative when research wages, product demand, and financing costs also change. Specify incumbent own-product innovation separately from acquisition of entrants or innovation in unrelated lines.
Use instruments shifting entrant innovation opportunities while satisfying stated exclusions for incumbent demand and research productivity; provide bounds if those exclusions are only partly defensible. Match incumbent research, product-quality changes, and entrant hazard, rather than inferring strategic effects from aggregate entry alone. Characterize mechanisms under which replacement risk discourages research or an escape-from-competition incentive raises it, and test whether their quantitative balance varies across technological distance. The target is a causal strategic-response function, not merely an accounting decomposition of growth into incumbent and entrant contributions.

For \(V\) twice continuously differentiable, an interior optimum satisfying \(V_r=0\) and \(V_{rr}<0\) has \(\partial r^*/\partial\lambda=-V_{r\lambda}/V_{rr}\). Its sign is therefore the sign of the cross-partial \(V_{r\lambda}\); this calculus identity is not the unresolved economic question. The sign depends on primitives and empirical causal responses and need not be universal across technologies. At a boundary or with multiple best responses, define finite differences or monotonicity of the best-response correspondence instead. General-equilibrium responses additionally include \(r_w^*dw/d\lambda+r_p^*dp/d\lambda\) and other endogenous prices. Entrant hazard should be measured as the arrival probability/intensity of innovations that displace the incumbent's own product, not all firm entry.
\subsection*{Variants and assumption changes}
The following are \emph{editorial variants}, not additional published open conjectures.
\begin{enumerate}
\item \emph{Replacement-only benchmark.} Let \(V(r,\lambda)=\pi(r)/(\rho+\lambda)-c(r)\) with \(\rho>0\), \(\pi'>0\), and a strictly concave objective. Then \(V_{r\lambda}<0\), yielding substitution locally; this is a solved conditional mechanism to compare with data.
\item \emph{Escape and endogenous wages.} Let entrant pressure change profits conditional on an incumbent's quality lead, and determine research wages in equilibrium. Identify the partial response at fixed wages separately from the total response, allowing opposite signs.
\end{enumerate}
\subsection*{References and source audit}
\begin{itemize}
\item Peter J. Klenow; Huiyu Li. \emph{Innovative Growth Accounting}. 2021. Locator: SectionVII, Conclusion, penultimate research paragraph on innovation arrival rates; published May 2021. Primary text checked. \url{https://www.journals.uchicago.edu/doi/full/10.1086/712325}.
\end{itemize}
Complementarity of creative destruction and incumbent innovation is explicit; the derivative design is editorial.
\begin{enumerate}\raggedright
\item\hypertarget{definitions:source:30007547:1}{} Peter J. Klenow; Huiyu Li. 2021. \emph{Innovative Growth Accounting}. SectionVII, Conclusion, penultimate research paragraph on innovation arrival rates; published May2021.
\url{https://www.journals.uchicago.edu/doi/full/10.1086/712325}.
DOI: \url{https://doi.org/10.1086/712325}.
\end{enumerate}

\subsection*{Common conventions: Source mathematical conventions}

These conventions apply to each entry unless explicitly replaced.

\textbf{Models and identification.} A primitive model \(m\) specifies all probability laws, feasible choices, information, equilibrium and measurement rules used by its target. The admissible class \(\mathcal M\) is fixed independently of outcomes. If \(P_O(m)\) is its observable probability law and \(\tau(m)\) the target, its identified set at observable law \(P\) is
\[
 \mathcal I_\tau(P)=\{\tau(m):m\in\mathcal M,\ P_O(m)=P\}.
\]
Point identification means that this set is a singleton. An empty set signals incompatibility of data and assumptions. Coordinate bounds need not describe the joint set. Bounds are sharp only if every retained target is attainable or arbitrarily approachable by compatible models. Finite bounds require explicit restrictions on unobserved quantities; unrestricted confounding, hidden wealth, extrapolation or arbitrary equilibrium selection can yield uninformative sets.

\textbf{Causal interventions.} A potential outcome \(Y_i(p)\) is the outcome under a complete policy regime \(p\), including timing, financing, information and market spillovers. The target population and units are fixed before treatment. With interference, use the complete assignment vector or a justified exposure mapping; individual no-interference notation alone cannot identify an economy-wide intervention. A difference of mean potential outcomes does not require the cross-world joint distribution. Conditioning on post-treatment migration, survival, enrollment or employment changes the estimand and needs separate justification.

\textbf{Statistical statements.} An estimator, confidence set, forecast or test specifies a sampling law, model class and loss/coverage criterion. Uniform \((1-\alpha)\)-coverage means \(\inf_{m\in\mathcal M}\Pr_m\{\tau(m)\in C_n\}\geq1-\alpha\), or an explicitly asymptotic version. The whole parameter space trivially has coverage, so informativeness requires a stated width, risk or power objective. A forecasting improvement is relative to a named benchmark, information set and evaluation population.

\textbf{Equilibrium and optimization.} An equilibrium comprises feasible plans, optimality, beliefs where needed, budget constraints and all market/resource clearing. The primitives or a stated benchmark define each correspondence \(\mathcal E(p;m)\). Nonemptiness is assumed or proved, never inferred just from compact policy sets. A selected-equilibrium target uses a declared selection rule; a robust target quantifies over all permitted equilibria. Use suprema or infima unless attainment is established. An \(\varepsilon\)-optimal policy is feasible and within \(\varepsilon\) of the finite optimum value. Report an empty feasible set separately.

\textbf{Welfare and accounting.} Normative utility, interpersonal normalization, social weights, numeraire, discounting and the use of public receipts are primitives. Behavioral utility need not coincide with the welfare criterion. Household welfare includes ownership income and rents once; adding profits again to compensating variation generally double-counts them. A monetary equivalent is defined by an explicit transfer and price convention, with monotonicity and range conditions. Average effects, mechanism contributions, transfers and welfare losses are different targets. Infinite sums require convergence and dynamic feasible plans require terminal or no-Ponzi conditions.

\textbf{Source versions.} A working-paper date, revision date and journal publication year are distinguished. A DOI for a journal version is not automatically the DOI of an earlier manuscript. Locators refer to the stated version and printed pages where specified. A metadata-only check does not verify a claim about a full-text conclusion. Every entry's source subsection narrows the provenance claim accordingly.
% END ECONOMICS STATEMENT
\end{document}
